\section*{Abstract}

In-context segmentation predicts the mask of a target structure from a few
annotated examples, so a single model can segment new classes without
retraining. For 3D medical images, compute becomes a central constraint,
since the cost of dense attention between target and context grows
quadratically with the number of voxels. This thesis develops an efficient 3D
in-context segmentation model around three questions: how to combine image
and label information, where to spend compute, and how to obtain training
tasks beyond the annotated classes.

First, we propose dual attention, which keeps image and label tokens in
separate streams that interact in every layer. At matched compute and with
fewer parameters, it outperforms merging the two once on 12 of 13 evaluation
splits, with the largest gains on new targets. Second, we propose a
coarse-to-fine cascade in which each level processes a single crop placed by
the previous level's prediction, which also enters the target's label tokens
as a prior. Training with the model's own predictions as prior outperforms
training with perturbed ground-truth masks. Third, we propose a synthetic
task generator that paints procedural shapes onto real scans; compared with a
fully synthetic canvas, it improves three of four lesion datasets.

We evaluate the model on held-out datasets and
compare it with Medverse. With matched refinement steps, our final model is
15 times faster and more accurate on average over the non-lesion datasets
(Dice 0.615 vs.\ 0.552), while Medverse remains more accurate on two of them.
Lesions remain difficult for both models and are the main open problem.